\begin{theorem}
\label{theorem-uncountable-nullstellensatz}
Let $k$ be a field. Let $S$ be a $k$-algebra generated over $k$
by the elements $\{x_i\}_{i \in I}$. Assume the cardinality of $I$
is smaller than the cardinality of $k$. Then
\begin{enumerate}
\item for all maximal ideals $\mathfrak m \subset S$ the field
extension $\kappa(\mathfrak m)/k$
is algebraic, and
\item $S$ is a Jacobson ring.
\end{enumerate}
If $I$ is finite, both conclusions hold for every field $k$, without
the cardinality restriction on the finite generating set.
\end{theorem}

\begin{proof}
For finite $I$, both assertions follow from the Hilbert
Nullstellensatz (Theorem \ref{theorem-nullstellensatz}) over every
field, regardless of the size of that finite set. Now assume $I$
is infinite and put $\kappa=\#I<\#k$.
Let $P=k[\{X_i\}_{i\in I}]$ be the polynomial ring mapping onto
the original algebra $S$ by $X_i\mapsto x_i$. Every monomial has
finite support and is specified by a finite list of pairs
$(i,e)$ with $i\in I$, $e\in\mathbf Z_{\geq0}$.
For each finite length $r$, the set of such lists has cardinality
at most $\kappa$, since $\kappa$ is infinite; the union over all
$r\geq0$ still has cardinality $\kappa$.
Distinct variables supply $\kappa$ distinct monomials, so the
monomial basis of $P$ has cardinality exactly $\kappa$.
The images of these monomials span $S$ and every quotient $S/\mathfrak m$.
In particular, for a maximal ideal $\mathfrak m\subset S$, its
residue field is exactly $S/\mathfrak m$ and has $k$-dimension
at most $\kappa$. These are maps from the same original polynomial
ring and algebra, without replacing $S$ by $P$.

\medskip\noindent
To arrive at a contradiction pick $T \in \kappa(\mathfrak m)$ transcendental
over $k$. Note that the $k$-linear map
$T : \kappa(\mathfrak m) \to \kappa(\mathfrak m)$
given by multiplication by $T$ has the property that $P(T)$ is invertible
for all monic polynomials $P(t) \in k[t]$.
Also, $\kappa(\mathfrak m)$ has dimension at most the cardinality of $I$
over $k$ since it is a quotient of the vector space
$k[\{x_i\}_{i \in I}]$ over $k$ (whose dimension is $\# I$ as we saw above).
This is impossible by Lemma \ref{lemma-dimension}.

\medskip\noindent
To show that $S$ is Jacobson we argue as follows. If not then
there exists a prime $\mathfrak q \subset S$ and an element $f \in S$,
$f \not \in \mathfrak q$ such that $\mathfrak q$ is not maximal
and $(S/\mathfrak q)_f$ is a field, see
Lemma \ref{lemma-characterize-jacobson}.
The field $L=(S/\mathfrak q)_f$ is generated over $k$ by the
images of the original $x_i$ and the inverse of the image of $f$.
Thus it has a generating set of cardinality at most
$\kappa+1=\kappa<\#k$. Apply the first assertion already proved
to the maximal ideal $(0)$ of the field $L$: the extension $L/k$
is algebraic. The canonical inclusion
$S/\mathfrak q\subset L\subset\operatorname{Frac}(S/\mathfrak q)$
identifies $L$ with the entire fraction field, since a field containing
$S/\mathfrak q$ contains the inverse of every nonzero element of
that domain. Hence $L=\kappa(\mathfrak q)$ canonically.
The prime $\mathfrak q$ lies over the maximal ideal $(0)$ of $k$,
so Lemma \ref{lemma-finite-residue-extension-closed} makes
$\mathfrak q$ maximal in $S$, a contradiction.
\end{proof}